% !TeX program = xelatex
\documentclass[UTF8,a4paper,11pt,fontset=fandol]{ctexart}
\IfFontExistsTF{SimSun}{\setCJKmainfont[AutoFakeBold=2]{SimSun}}{}
\IfFontExistsTF{Microsoft YaHei}{\setCJKsansfont{Microsoft YaHei}}{}
\usepackage[margin=22mm,headheight=15pt]{geometry}
\usepackage{hyperref,xcolor,fancyhdr,needspace,enumitem}
\definecolor{schoolgreen}{HTML}{126831}
\hypersetup{colorlinks=true,urlcolor=schoolgreen,pdftitle={AI伴学第一课教师讲稿},pdfauthor={华光辉}}
\pagestyle{fancy}\fancyhf{}\fancyhead[L]{\small AI伴学第一课·教师讲稿}\fancyhead[R]{\small 2026.09.12}\fancyfoot[C]{\thepage}
\ctexset{section={format=\large\bfseries\color{schoolgreen},beforeskip=1.2ex,afterskip=.6ex}}
\setlength{\parindent}{2em}\setlength{\parskip}{.4em}\linespread{1.22}
\newcommand{\tip}[1]{\par{\small\color{black!65}\noindent\textbf{操作与备课：}#1\par}}
\begin{document}
\begin{center}{\LARGE\bfseries\color{schoolgreen}AI伴学第一课}\par\vspace{4mm}{\large 教师讲稿}\par\vspace{3mm}华光辉\quad 40分钟\end{center}






这份讲稿对应同目录29页Beamer，研究信息核查至2026年9月11日。正文可直接讲述，操作与备课提示供教师查看。可按自己的口头习惯改词；系统演示中的“这里”“这一步”，要对应课前选好的真实资料。

40分钟路线：课程内容与要求7分钟；AI历史、神经网络入门和两段完整视频17分钟；IMO成绩引入1分钟；NS方程通俗介绍4分钟；Scholar演示10分钟；结束书写1分钟。第21页讲完直接跳第29页。第22—28页是集合题拓展，不计入这40分钟。

课前准备：提前打开两段B站视频并检查声音；启动Scholar，打开已登录的本地实例；预选一页资料和一条非敏感练习记录；另备已保存的对话，以防网络响应较慢。视频是联网链接，不是PDF内嵌文件。

本版删去2021、2023年的数学成果，以及单位距离、雅可比猜想的讲解。NS部分只讲水流情境、一个高中函数和AI参与的研究，不展示微分方程与证明细节。

Scholar的10分钟包含真实操作、对照原文、学生复述或追问，无需持续读稿。到39分钟直接回第21页收尾；集合题通常留到下次课。

第3—4页为教学大纲，依据《教学内容（暂定）.docx》的18课时主表归纳。保留原有主题、先后顺序和课时范围，以学生能听懂的语言介绍；文档补充页的旧第一课时间表不作为本次排课指令。


\clearpage

\par\noindent\begin{minipage}{\linewidth}

\section*{第1页　AI伴学 第一课}

{\small\noindent 00:00—00:30\par}

同学们好，这门选修课叫“AI伴学”。接下来一个学期，我们会用AI读一些东西、讨论一些问题，也会做数学。今天我先把这门课怎么上、期末怎么评价说清楚，再带大家看看AI最近做到了什么。最后，我演示一下自己使用的学习系统。

\tip{停留封面，先和学生建立正常的课堂交流，不急着报工具名称。}

\end{minipage}\par\vspace{2mm}

\par\noindent\begin{minipage}{\linewidth}

\section*{第2页　这门课怎么上}

{\small\noindent 00:30—01:00\par}

今天主要就是这三件事。大家不需要先会编程，也不用现在就听懂前沿论文。听的时候可以留意：有哪些问题你想继续了解？后面我会给大家看一道集合题；如果时间不够，我们下次接着讲。

\tip{不逐字读目录。集合题明确为机动内容。}

\end{minipage}\par\vspace{2mm}

\par\noindent\begin{minipage}{\linewidth}

\section*{第3页　这学期学什么（一）}

{\small\noindent 01:00—02:00\par}

先看看这学期具体学什么。前两课，我们认识AI，再配置一个适合自己的学习助手，学会怎样提问、要提示和检查回答。

接下来有三个数学主题。复数从x平方加1等于0开始，再把它画到平面上，看看旋转和伸缩；图论从一笔画开始，把路线画成点和边；迭代与分形，则是反复代入公式、重复画图，观察出现了什么规律。

这些内容不要求大家提前学过。我们会从老师示范，逐步过渡到一起制定学习步骤，再由大家自己提出问题、安排探索。

\tip{依据《教学内容（暂定）.docx》主表第1—10课时归纳。用约1分钟介绍主题，不现场解释复数、一笔画条件或分形定义。前两课内容仍按当前课件及后续配置课执行，不照搬文档补充页的旧第一课时间表。}

\end{minipage}\par\vspace{2mm}

\par\noindent\begin{minipage}{\linewidth}

\section*{第4页　这学期学什么（二）}

{\small\noindent 02:00—03:00\par}

后半学期，我们把前面练习的方法用回课内。第11课会选已经学过的集合、函数等内容，检查自己是不是真的理解。第12到15课，完整做一轮：先独立作答，再分析错因、做针对性的练习，隔几天不用AI再做一道新题。

第16课练习条件转化、分类讨论和规范表达，范围只用当时已经学过的知识。最后两课，大家展示自己的学习过程，完成独立任务，再为一个想学的新主题设计入门计划。

目前暂定18课，每课40分钟，具体课内题目会跟着我们的学习进度调整。下面再说期末怎样评价。

\tip{依据主表第11—18课时归纳。第17课侧重方法成果展示，第18课包含分层独立任务与新主题计划；新主题计划不等于已经掌握。两页介绍与后面的40/30/30评价衔接，不新增考试范围或提交要求。}

\end{minipage}\par\vspace{2mm}

\par\noindent\begin{minipage}{\linewidth}

\section*{第5页　期末评价与学分}

{\small\noindent 03:00—05:30\par}

先说大家最关心的期末评价。按照现在的课程安排，学习档案占40\%，独立数学任务占30\%，探究与迁移展示占30\%。

学习档案，是让别人看得出你怎么学。比如一道题，第一次写到哪里，后来问了什么，为什么改了某一步，这些都可以保存。独立任务，是检查你不借助AI时能完成什么。展示的时候，可以选一个自己认真研究过的问题，把它讲明白，并说明这个方法还能用在哪里。

所以，平时就要留下记录，不必等到期末再回忆。学分具体是多少、怎样认定，最终按学校公布的要求执行，申报材料里这一项还没有核定。我不会今天先给大家承诺一个分数或学分数。

\tip{可留20—30秒回答评价问题。没有材料依据的合格线、考勤比例、提交截止日期不临时编造。}

\end{minipage}\par\vspace{2mm}

\par\noindent\begin{minipage}{\linewidth}

\section*{第6页　学习记录可以这样写}

{\small\noindent 05:30—07:00\par}

一份学习记录可以很简单。比如你做一道集合题，写到一半发现自己不知道有没有漏掉情况，就先把这一步存下来。然后问AI：我这样分类完整吗？它提出建议以后，你要看建议有没有道理，再自己补上。

隔一天，不看那段对话，再做一次。这时你就知道，自己是看懂了一次回答，还是已经能够独立解决。记录里允许有错，也应该保留修改。否则到期末，你自己也很难说清楚是怎样学会的。

\tip{指着“尝试”和“回头再做”两行讲，不扩展成额外的评分规定。}

\end{minipage}\par\vspace{2mm}

\par\vspace{4mm}

\par\noindent\begin{minipage}{\linewidth}

\section*{第7页　人工智能并不是这几年才出现的}

{\small\noindent 07:00—07:45\par}

人工智能并不是这几年才出现的。1956年的达特茅斯夏季研究计划，通常被看作一个重要起点。研究者已经在问：机器能不能使用语言、形成概念、解决问题？早期一种想法，是把知识和规则写给计算机。但认出照片里的猫，很难只靠几条简单规则。另一条思路，就是让计算机从例子中学习。

\tip{“1956年”指夏季研究计划；提案写于1955年。不要讲成AI在某一天突然被发明。}

\end{minipage}\par\vspace{2mm}

\par\noindent\begin{minipage}{\linewidth}

\section*{第8页　让计算机从例子中学习}

{\small\noindent 07:45—08:15\par}

2012年的AlexNet在图像识别中取得重要进展，2016年的AlphaGo以4比1战胜李世石。这些系统都用到了神经网络。这个名字听起来复杂，我们先看其中一次计算，其实只需要高中学过的代数。

\tip{4比1是2016年与李世石五番棋的总比分。不要把AlphaGo说成单纯穷举全部棋局。}

\end{minipage}\par\vspace{2mm}

\par\noindent\begin{minipage}{\linewidth}

\section*{第9页　一个“神经元”怎样计算}

{\small\noindent 08:15—09:15\par}

图片可以用一串数字表示。我们先只取两个输入，看看这个小单元怎么处理。每个数先乘一个系数，再相加，最后加一个常数。这和我们学过的代数式很像，系数叫权重，最后的常数叫偏置。

代入1和0.5，算出来是0.3。接着过一道简单规则：负数变成0，非负数保留。这叫激活函数，这里用的是ReLU。如果刚才算出来是负0.2，输出是多少？对，是0。这个数还要传给后面的计算。

\tip{指着箭头和算式讲，留5秒让学生算负数的情况。示例参数由教师给定，0.3不是识别概率。无需讲导数或反向传播公式。}

\end{minipage}\par\vspace{2mm}

\par\noindent\begin{minipage}{\linewidth}

\section*{第10页　连成网络，再用例子训练}

{\small\noindent 09:15—10:15\par}

把许多这样的小单元连起来，上一层算出的结果交给下一层，就组成一个网络。中间可以有很多层。激活函数让这些计算能够表达比一条直线更复杂的关系。

训练时，先给它一个例子，让它算，再和已知答案比较。结果差得多，就根据误差调整权重和偏置，使用许多例子反复练习。训练完，我们还要拿它没见过的例子来检查效果。就像做过的题会了，还要看换一道题能不能做。下面回到AlphaGo，看看人们怎样看待它的表现。

\tip{这里讲的是监督学习的简化情形。输入节点只传入数据；隐藏层做计算，输出层按任务设计。参数不是教师逐个手调，而是由训练算法更新；每一步不保证所有例子都改善。学生追问时再介绍“反向传播计算参数怎样影响误差”，不展开公式。}

\end{minipage}\par\vspace{2mm}

\par\noindent\begin{minipage}{\linewidth}

\section*{第11页　纪录片《AlphaGo》预告}

{\small\noindent 10:15—12:00（含视频1分30秒）\par}

我们先看这段有中文字幕的预告。看完后，大家想一想：一台机器下出人没想到的棋，对棋手来说意味着什么？

\tip{点击“播放视频”，全屏播放，结束后回到PDF。15秒左右用于引入与切换，不预设学生会怎样回答。}

\end{minipage}\par\vspace{2mm}

\par\noindent\begin{minipage}{\linewidth}

\section*{第12页　Hinton谈自己的研究经历}

{\small\noindent 12:00—22:45（含视频10分27秒）\par}

下面听Hinton谈自己的研究经历。今天熟悉的技术，研究过程中也经历过曲折。听一听他怎样回顾当时的选择。

\tip{完整播放10分27秒。课前打开网页并确认音量；视频内容未逐段转写，不在讲稿里虚构某句台词或片段时间戳。若讨论视频标题中的历史假设，要说“这是对另一种历史可能性的设想”。}

\end{minipage}\par\vspace{2mm}

\par\noindent\begin{minipage}{\linewidth}

\section*{第13页　今天的AI为什么能和我们对话}

{\small\noindent 22:45—24:00\par}

刚才看到的神经网络，只是一个很小的示意。今天的聊天式AI，使用的网络结构复杂得多。2017年提出的Transformer，用注意力机制处理上下文中信息之间的联系。模型训练后，可以接着上下文生成文字，也能写代码和解释问题。

再接上检索、计算和证明检查工具，就能边尝试边检查。不过，回答像标准答案，还不能代替每一步都有理由。接下来我们看看，它们在数学研究里做出了哪些进展。

\tip{可等待一个简短回答，控制在15秒内。不把语言模型说成拥有保证正确的内置知识库。}

\end{minipage}\par\vspace{2mm}

\par\noindent\begin{minipage}{\linewidth}

\section*{第14页　从解竞赛题，到研究新问题}

{\small\noindent 24:00—25:00\par}

2025年，增强版Gemini Deep Think做国际数学奥林匹克试题，六道解出了五道，得了35分，达到当年的金牌标准。大家要注意，竞赛题很难，但它已经有答案和评分标准。研究问题更进一步，可能连数学家也还不知道答案。

下面我们不看复杂的代数问题，就从大家熟悉的水流，说起一个研究问题。

\tip{R13。模型达到金牌标准，不是作为学生参赛获得奖牌。这一页只作过渡，不展开竞赛解法。}

\end{minipage}\par\vspace{2mm}

\par\noindent\begin{minipage}{\linewidth}

\section*{第15页　NS方程：用数学描述水怎样流动}

{\small\noindent 25:00—26:15\par}

想象用勺子搅动一杯水。水会转起来，停下以后，又会慢慢平静下来。要描述这个过程，不能只说整杯水朝哪个方向走，因为杯子不同位置的水，速度和方向都可能不同。

NS方程的基本想法，和物理课里的牛顿第二定律相通：研究一小部分水受到什么作用、运动怎样改变。压力差会推挤它，邻近的水会因为黏性彼此牵制，也可能有重力或其他外力。方程把这些作用联系起来，描述各处的水接下来怎样流动。

这里不需要记方程，只要知道：它是在用数学追踪流体的运动。

\tip{R8、R10。用手势表示不同位置的流速即可，不做真实流体实验。杯中水是引入情境，不是论文研究区域；黏性不是所有流速都单调减小的保证。不引入偏导数、散度或拉普拉斯算子。}

\end{minipage}\par\vspace{2mm}

\par\noindent\begin{minipage}{\linewidth}

\section*{第16页　难在哪里：一直算下去，会发生什么}

{\small\noindent 26:15—27:45\par}

难题在于：从一个很光滑的状态开始，按方程一直算下去，流动能始终保持光滑吗？还是可能在有限时间内出现奇异性？这里说的是速度等量变化得是否平顺，不是水面有没有波纹。我们先用一个高中函数，理解其中一种“失控”的意思。

看表格，t取0.9，y是10；取0.99，y是100；再接近1一点，就变成1000。请大家接着算：t取0.9999，y是多少？对，是10000。时间没有走到无穷远，只是在接近1，函数值却可以超过任何预先给定的数。

数学上的“有限时间爆破”，指的就是这种无界增长。这不是水真的爆炸，也不是说只要水流很乱就发生了爆破。这个函数只是帮助我们理解词义，真实的NS问题复杂得多。

\tip{留约10秒让学生口算。正确答案10000。只讨论t从小于1的一侧趋近1；t=1不在定义域。光滑性问题比单个速度值是否有界更一般，这里用速度爆破介绍新论文涉及的一种情形。湍流或普通漩涡不等于数学奇异性。}

\end{minipage}\par\vspace{2mm}

\par\noindent\begin{minipage}{\linewidth}

\section*{第17页　AI参与了什么}

{\small\noindent 27:45—29:00\par}

今年9月8日，OpenAI公布了一项研究，报告AI系统构造了一种特殊的三维流动：在光滑外力作用下，数学模型中的速度能够在有限时间内无界增长。团队同时公开了证明论文，以及供计算机逐步核查推理的代码。

这件事让我们看到，AI已经在参与寻找新构造和证明。但这是刚公布的结果，仍需要进一步独立审查。现实中的水也不会真的达到无穷大的速度；数学模型出现这样的结果，提醒我们要分清模型和现实。

今天不要求大家掌握证明。知道这个问题在问什么、AI参与了什么，就够了。下面回到我们自己的学习，看看我怎样使用Scholar。

\tip{R8—R10。2026-09-11复查作者公布页和Clay页面；后者仍列Active。精确条件是三维不可压、有适当光滑外力、初始静止且动能有界。课堂不讲C、D表述、能量估计或Lean语法；没有独立复核论文和代码，不宣布正式获奖或已获普遍认可。}

\end{minipage}\par\vspace{2mm}

\par\noindent\begin{minipage}{\linewidth}

\section*{第18页　我怎样使用Scholar学习}

{\small\noindent 29:00—30:00\par}

接下来介绍我自己搭建的Scholar系统。前面讲的是数学研究，现在看日常学习：我把教材、笔记和学习记录放在一起。读到不懂的地方就问，做题卡住了就先保存尝试，再要一点提示，过几天还能回来看。

我今天只演示一个小过程。大家先看我怎样问、怎样看回答，再想一想，换成你正在学的内容，你会在哪一步需要帮助。

\tip{切换到教师已有的本地实例。使用真实的非敏感资料；不展示账号密钥、学生姓名和成绩。不臆造系统已加载的书名。}

\end{minipage}\par\vspace{2mm}

\par\noindent\begin{minipage}{\linewidth}

\section*{第19页　演示：读到一句没看懂的话}

{\small\noindent 30:00—34:00\par}

这里我打开一份准备好的资料。我们只看这一句话，不把整本书的问题一次丢给AI。我先说自己的理解，再问这一步用了什么条件。

现在看它的回答。它指出的内容能不能在原文里找到？如果举了例子，这个例子满足原来的条件吗？有一处没懂，我就继续追问这一处。

读书时，Ask \& Understand适合这样连续讨论。看懂回答以后，我还会用自己的话再说一遍，看看自己到底能不能解释。

\tip{Library$\to$Open reader$\to$定位预选页面$\to$选中一句话$\to$Explain/Check me或输入课件提示。约2分钟真实操作与原文对照，留1分钟请学生用自己的话复述或提出追问。响应超过30秒，打开课前已有记录并说明这是以前保存的对话。}

\end{minipage}\par\vspace{2mm}

\par\noindent\begin{minipage}{\linewidth}

\section*{第20页　演示：一道题暂时做不下去}

{\small\noindent 34:00—39:00\par}

再看练习的情况。我先写出已经想到的部分，保存下来。这时如果确实做不下去，就点一次“给我一个提示”。

拿到提示后，我先自己接着写，而不是连续点到完整答案出现。改动了哪一步，也要留下来。这样回头看记录，才能知道自己原来为什么卡住。

最后安排一下复习。系统有1天后和7天后的入口；时间到了，我重新做一次。是不是学会了，要看那时自己能否完成。大家以后可以用同样的方法处理自己的问题，不一定和我演示的题一样。

\tip{Practice \& Prove$\to$Attempt$\to$Save attempt$\to$Request one hint$\to$保存修改$\to$Schedule 1d + 7d reviews$\to$Reports。约3分钟真实操作，留1分钟请学生判断提示是否解决了卡点。到39分钟必须回第21页；只有提前完成，才点集合题。}

\end{minipage}\par\vspace{2mm}

\par\noindent\begin{minipage}{\linewidth}

\section*{第21页　下课前}

{\small\noindent 39:00—40:00\par}

今天就先到这里。请大家写下一个这学期想弄懂的问题。可以是你做不下去的一道题，也可以是读书时没看懂的一句话，或者今天哪一个数学问题让你想继续了解。

不用写得很大，具体一点最好。下次课，我们带着自己的问题，配置和试用学习助手。

\tip{给学生约30秒书写。点击“结束”跳到第29页，避免顺序进入集合附录。}

\end{minipage}\par\vspace{2mm}

\par\vspace{4mm}

\par\noindent\begin{minipage}{\linewidth}

\section*{第22页　选讲：集合中的一个难题}

{\small\noindent 机动；首次展示约1分钟\par}

如果还有几分钟，我们看这道集合题。先别急着看第三问。这里T(B)表示B里面所有元素的和。第一问是列出所有可能的集合；第二问要求选三个元素；第三问才是一般性的证明。

注意，集合里的元素互不相同，B还不能是空集。让AI读题时，也要先把这些条件读准确。

\tip{优先只展示22—23页，完整解答供下次课。题面由原印刷照片整理，照片本身不放进课件。}

\end{minipage}\par\vspace{2mm}

\par\noindent\begin{minipage}{\linewidth}

\section*{第23页　选讲：先让AI帮哪一步}

{\small\noindent 机动；约1分钟\par}

我不会一开始就让它把整道题讲完。先问它有没有读对题，再问第一、二问怎样分类。到第三问，可以要求只给一个入口。

假如它建议“把和为2n的数配对”，我们就检查这句话后面还有没有漏掉情况。AI给的思路，我们照样可以追问：你这一步为什么一定成立？

\tip{教学目的：发现量词、非空条件和分类中的遗漏。可把一次问答放进Scholar，不要求当场生成完整证明。}

\end{minipage}\par\vspace{2mm}

\par\noindent\begin{minipage}{\linewidth}

\section*{第24页　选讲：第一问的完整列举}

{\small\noindent 拓展；约1分钟\par}

按元素个数分类最稳妥。一个元素只有8；两个元素有1和7、2和6；三个元素有1、2、5。四个及以上为什么不用再找？最小的四个数加起来已经是14，超过8了。

这样不仅列出了四个集合，还说明了为什么没有遗漏。

\tip{目标：按元素个数分类。易错：漏掉单元素集合或三元素集合。检查学生能否解释“四个及以上不可能”。}

\end{minipage}\par\vspace{2mm}

\par\noindent\begin{minipage}{\linewidth}

\section*{第25页　选讲：第二问的完整列举}

{\small\noindent 拓展；约1分钟\par}

这次固定取三个元素。最小和是12，最大和是30，所以只用检查12和24。和为12的有一组，和为24的有三组，一共四组。

同一组元素换个顺序，并不会变成一个新集合。

\tip{目标：用最小值、最大值缩小枚举范围。易错：重复排列计数。检查学生能否先说出只需考察的两个和值。}

\end{minipage}\par\vspace{2mm}

\par\noindent\begin{minipage}{\linewidth}

\section*{第26页　选讲：配对法的关键分支}

{\small\noindent 拓展；约2分钟\par}

把1到2n减1按和为2n配对，留下n单独一组。如果A里有一整对，就完成了。真正需要处理的是另一种情况：每一对最多取一个数。

这里总共有n组，而A恰好有n个元素，所以只能每组取一个，特别是n必须在A里。不能说“由抽屉原理，一定有一组取了两个”——n个元素放进n组，正好每组一个也完全可能。

\tip{目标：完整分类与抽屉原理的条件。难点：第二种情况。学生若指出错误的抽屉原理应用，就达到了这一页的判断目标。}

\end{minipage}\par\vspace{2mm}

\par\noindent\begin{minipage}{\linewidth}

\section*{第27页　选讲：一个余数引理}

{\small\noindent 拓展；约3分钟\par}

下面用一个引理。假如n减1个整数，无论选哪几个相加，都不能被n整除，那么它们除以n的余数必须一样。

把它们随意排好，依次取前一个、前两个、前n减1个数的和。这些部分和的余数不能是0，也不能重复。否则两个部分和相减，就会得到一组和能被n整除的数。于是它们恰好占满所有非零余数。

现在交换前两个数，后面的部分和都不变，只有第一个改变。因为仍要占满同一组余数，第一个余数也不能变。于是这两个数同余。任何两个数都可以排到前两位，引理就成立了。

\tip{这里使用n$\ge$4，项数至少为3。关键是“任意排列”与“只有第一个部分和改变”，不是只对某个固定顺序观察。}

\end{minipage}\par\vspace{2mm}

\par\noindent\begin{minipage}{\linewidth}

\section*{第28页　选讲：完成第三问的证明}

{\small\noindent 拓展；约2分钟\par}

回到第二种情况，n已经在A中。假设没有我们要的B，把n取走，剩下的集合记为C。C里也不能有一组数的和是n的倍数：如果这个倍数是偶数，它本身就合要求；如果是奇数，再加上n，也合要求。

所以C满足刚才的引理，所有元素模n同余。可是从1到2n减1，同一个非零余数最多只有两个数；C有n减1个互不相同的元素，至少有三个，矛盾。

注意，这里才真正用到了n至少为4。证明到最后，我们应该能指出每个条件用在了哪里。

\tip{目标：反证法、奇偶分类和引理应用。第三问难度较高，首次课不要求学生当场掌握；可作为后续分析AI解答是否漏条件的材料。}

\end{minipage}\par\vspace{2mm}

\par\noindent\begin{minipage}{\linewidth}

\section*{第29页　结束页}

{\small\noindent 结束\par}

下次课见。

\tip{正常40分钟路线由第21页直接跳转至此。}

\end{minipage}\par\vspace{2mm}

\par\vspace{5mm}\section*{备课文献}

以下链接供教师备课，不在课件中放映。研究结论的精确范围见同目录《前沿数学核查记录.md》。

\begin{enumerate}[label={},leftmargin=0pt,itemsep=6pt]

\item R8　\href{https://openai.com/index/navier-stokes-solution/}{NS研究公布说明}

\item R9　\href{https://github.com/openai/NavierStokesAndEuler}{作者的Lean代码与构建说明}

\item R10　\href{https://www.claymath.org/millennium/navier-stokes-equation/}{Clay的NS问题页面}

\item R13　\href{https://deepmind.google/blog/advanced-version-of-gemini-with-deep-think-officially-achieves-gold-medal-standard-at-the-international-mathematical-olympiad/}{2025年IMO金牌标准}

\item R14　\href{https://developers.google.com/machine-learning/crash-course/neural-networks/nodes-hidden-layers}{神经网络：节点、权重与隐藏层}

\item R15　\href{https://developers.google.com/machine-learning/crash-course/neural-networks/activation-functions}{神经网络：激活函数}

\end{enumerate}\end{document}